\documentclass[12pt,a4paper]{article}
\usepackage[utf8]{vietnam}
\usepackage{amsmath,amssymb}
\usepackage{tabularx}
\usepackage{tikz}
\usetikzlibrary{calc,angles,quotes,arrows.meta,patterns}
\usepackage{xcolor}
\usepackage{enumitem}
\usepackage{fancyhdr}
\usepackage{graphicx}
\usepackage[top=0.9cm, bottom=0.9cm, left=1.2cm, right=1.2cm]{geometry}
\usepackage[dethi]{ex_test}
\SAOPTN{widthSA=3.2,heightSA=0.62}

\setlength{\headheight}{15pt}
\setlength{\parskip}{0pt}

\pagestyle{fancy}
\fancyhf{}
\rhead{\textit{Môn Toán 9 -- Đề số 02}}
\lhead{\textit{Bộ đề ôn tập Giữa học kì I}}
\cfoot{\thepage}

\begin{document}

% --- KHUNG TIÊU ĐỀ ĐỀ THI ---
\noindent
\begin{tabularx}{\linewidth}{X r}
\textbf{TRƯỜNG THCS: \dotfill} & \textbf{ĐỀ KIỂM TRA ĐỊNH KỲ GIỮA HỌC KỲ I} \\
\textbf{Họ và tên: \dotfill} & \textbf{NĂM HỌC 2026 -- 2027} \\
\textbf{Lớp: \dotfill \quad SBD: \dotfill} & \textbf{Môn: TOÁN 9} \\
\hline
\end{tabularx}
\vspace{2mm}

\begin{center}
  \textbf{\large ĐỀ ÔN TẬP GIỮA HỌC KỲ I -- ĐỀ SỐ 02}\\
  \textit{Thời gian làm bài: 90 phút (Không kể thời gian phát đề)}
\end{center}
\vspace{1mm}

% --- PHẦN I ---
\noindent\textbf{\large PHẦN I. CÂU TRẮC NGHIỆM NHIỀU PHƯƠNG ÁN LỰA CHỌN (3,0 điểm)}

\textit{Thí sinh trả lời từ câu 1 đến câu 12. Mỗi câu hỏi thí sinh chỉ chọn một phương án duy nhất.}

\vspace{1mm}
\begin{ex}
Trong các phương trình sau, phương trình nào không phải là phương trình bậc nhất hai ẩn?
\choice
{$2x - 3y = 5$}
{$0x + 2y = 4$}
{$2x - 0y = 3$}
{\True $0x - 0y = 6$}
\end{ex}

\begin{ex}
Hệ phương trình nào dưới đây là hệ hai phương trình bậc nhất hai ẩn?
\choice
{$\begin{cases} 2x + 2y = 3 \\ x - 3y = 18 \end{cases}$ (với bậc khác)}
{$\begin{cases} 0{,}5x - 0{,}2y = -0{,}1 \\ 0x + 0y = -0{,}4 \end{cases}$}
{\True $\begin{cases} -4x + 7y = -10 \\ 3x + 8y = -19 \end{cases}$}
{$\begin{cases} 2x^2 + 3y = 2 \\ 31x + 5y = -1 \end{cases}$}
\end{ex}

\begin{ex}
Trong các hệ hai phương trình bậc nhất hai ẩn sau, hệ phương trình nào nhận cặp số $(-1; -2)$ là nghiệm?
\choice
{$\begin{cases} 12x - 3y = -6 \\ 5x = 5 \end{cases}$}
{$\begin{cases} 0{,}2x - 3y = 0{,}7 \\ 0{,}8x - 2y = 2 \end{cases}$}
{\True $\begin{cases} -x + y = -1 \\ 3x + y = -5 \end{cases}$}
{$\begin{cases} 3x + 2y = 31 \\ 5x + y = -1 \end{cases}$}
\end{ex}

\begin{ex}
Điều kiện xác định của phương trình $\frac{x}{x-2} - \frac{2}{x^2-4} = \frac{1}{x}$ là:
\choice
{$x \neq \pm 2$}
{$x \neq 2$}
{$x \neq 0$}
{\True $x \neq 0$ và $x \neq \pm 2$}
\end{ex}

\begin{ex}
Cho tam giác $ABC$ vuông tại $B$, khi đó $\sin C$ bằng:
\choice
{$\frac{AC}{BC}$}
{\True $\frac{AB}{AC}$}
{$\frac{AB}{BC}$}
{$\frac{BC}{AC}$}
\end{ex}

\begin{ex}
Cho $\alpha$ và $\beta$ là hai góc phụ nhau. Khẳng định nào sau đây đúng?
\choice
{$\sin \alpha = \tan \beta$}
{$\sin \alpha = \cot \beta$}
{$\cos \alpha = \cot \beta$}
{\True $\sin \alpha = \cos \beta$}
\end{ex}

\begin{ex}
Một ngân hàng đang áp dụng lãi suất gửi tiết kiệm kì hạn 12 tháng là $7{,}4\%/\text{năm}$. Bà Mai dự kiến gửi một khoản tiền và cần số tiền lãi hằng năm ít nhất là $60\text{ triệu đồng}$. Hỏi số tiền bà Mai cần gửi tiết kiệm ít nhất là bao nhiêu (làm tròn đến triệu đồng)?
\choice
{$810\text{ triệu đồng}$}
{$800\text{ triệu đồng}$}
{\True $811\text{ triệu đồng}$}
{$700\text{ triệu đồng}$}
\end{ex}

\begin{ex}
Biết rằng $m > n$ với $m, n$ bất kỳ, khẳng định nào sau đây đúng?
\choice
{\True $m - 3 > n - 3$}
{$m + 3 < n + 3$}
{$m - 2 < n - 2$}
{$n + 2 > m + 2$}
\end{ex}

\begin{ex}
Bất phương trình nào sau đây là bất phương trình bậc nhất một ẩn?
\choice
{$\frac{1}{2}x^2 - y < 0$}
{$10 < \frac{2}{y} - x$}
{\True $3y - 1 < 4$}
{$0y + 8 \ge 0$}
\end{ex}

\begin{ex}
Tỉ số lượng giác bằng $\tan 46^\circ$ là:
\choice
{$\tan 46^\circ$}
{$\cos 44^\circ$}
{$\tan 44^\circ$}
{\True $\cot 44^\circ$}
\end{ex}

\begin{ex}
Biểu thức $A = \sin 35^\circ + \cos 35^\circ - \sin 55^\circ - \cos 55^\circ$ có giá trị là:
\choice
{$\sin 40^\circ - \cos 40^\circ$}
{$-80^\circ$}
{\True $0$}
{Không tính được}
\end{ex}

\begin{ex}
Một cột điện cao $4\text{ m}$, tia nắng mặt trời chiếu xuống mặt đất tạo với mặt đất góc $35^\circ$. Chiều dài bóng của cột điện trên mặt đất (làm tròn đến hai chữ số thập phân) là:
\choice
{$4{,}40\text{ m}$}
{$3{,}41\text{ m}$}
{\True $5{,}71\text{ m}$}
{$2{,}09\text{ m}$}
\end{ex}


\vspace{1mm}
\noindent\textbf{BẢNG TRẢ LỜI TRẮC NGHIỆM PHẦN I:}
\begin{center}
\begin{tabular}{|c|c|c|c|c|c|c|c|c|c|c|c|c|}
\hline
\textbf{Câu} & \textbf{1} & \textbf{2} & \textbf{3} & \textbf{4} & \textbf{5} & \textbf{6} & \textbf{7} & \textbf{8} & \textbf{9} & \textbf{10} & \textbf{11} & \textbf{12} \\
\hline
\textbf{Chọn} & \rule{0pt}{0.65cm} & & & & & & & & & & & \\
\hline
\end{tabular}
\end{center}
\vspace{2mm}

\setcounter{ex}{0}
\noindent\textbf{\large PHẦN II. CÂU TRẮC NGHIỆM ĐÚNG SAI (2,0 điểm)}

\textit{Thí sinh trả lời câu 1 và câu 2. Trong mỗi ý a), b), c), d) ở mỗi câu, thí sinh chọn đúng hoặc sai.}

\vspace{1mm}
\begin{ex}
Giải hệ phương trình $\begin{cases} 2x - 6y = 5 \\ 2x - 6y = 4 \end{cases}$ bằng phương pháp cộng đại số. Xét tính Đúng/Sai của các khẳng định sau:
\choiceTF
{Nhân hai vế của phương trình thứ hai với 2, ta được $\begin{cases} 2x - 6y = 5 \\ 2x - 6y = 4. \end{cases}$}
{\True Trừ từng vế phương trình thứ nhất cho phương trình thứ hai của hệ, ta được $0x = 1$}
{Phương trình $0x = 1$ có vô số nghiệm}
{Nghiệm tổng quát của hệ phương trình đã cho là $(6y + 5; 2x - 4)$ với $x \in \mathbb{R}$}
\end{ex}

\begin{ex}
Cho bất phương trình $(5m - 2)x < 1$. Xét tính Đúng/Sai của các khẳng định sau:
\choiceTF
{Bất phương trình đã cho là bất phương trình bậc nhất ẩn $x$ với mọi $m \in \mathbb{R}$}
{\True Khi $m = 1$, bất phương trình đã cho có nghiệm là $x < \frac{1}{3}$}
{Khi $m = -1$, bất phương trình đã cho có nghiệm là $x < -\frac{1}{7}$}
{Khi $m = -2$, bất phương trình đã cho có nghiệm nguyên lớn nhất là $-1$}
\end{ex}


\noindent\textbf{BẢNG TRẢ LỜI PHẦN II:}
\begin{center}
\begin{tabular}{|c|c|c|c|c||c|c|c|c|c|}
\hline
\textbf{Câu 1} & \textbf{a)} & \textbf{b)} & \textbf{c)} & \textbf{d)} & \textbf{Câu 2} & \textbf{a)} & \textbf{b)} & \textbf{c)} & \textbf{d)} \\
\hline
\textbf{Đ/S} & \rule{0pt}{0.6cm} & & & & \textbf{Đ/S} & & & & \\
\hline
\end{tabular}
\end{center}
\vspace{2mm}

\setcounter{ex}{0}
\noindent\textbf{\large PHẦN III. CÂU HỎI TRẢ LỜI NGẮN (2,0 điểm)}

\textit{Thí sinh trả lời câu 1 và câu 2. Điền kết quả vào ô trống tương ứng.}

\vspace{1mm}
\begin{ex}
Cho hệ phương trình $\begin{cases} 2x - y = 0 \\ x + y = 6. \end{cases}$ Tìm nghiệm $(x; y)$ của hệ. \shortans{$(2; 4)$}
\end{ex}

\begin{ex}
Cho tam giác $ABC$ có $AB = 5\text{ cm}$, $BC = 12\text{ cm}$ và $CA = 13\text{ cm}$. Tính số đo góc $C$ (làm tròn đến độ). \shortans{$23^\circ$}
\end{ex}


\setcounter{ex}{0}
\noindent\textbf{\large PHẦN IV. TỰ LUẬN (3,0 điểm)}

\textit{Thí sinh trình bày chi tiết lời giải vào giấy làm bài.}

\vspace{1mm}
\begin{ex} (1,0 điểm)
Giải các phương trình sau:
\begin{enumerate}[label=\textbf{\alph*)}, leftmargin=0.8cm, itemsep=2pt]
\item (0,5 điểm) $2x(x - 3) + 5(x - 3) = 0$
\item (0,5 điểm) $\frac{x}{x + 2} + \frac{3}{x - 2} = \frac{3x + 4}{x^2 - 4}$
\end{enumerate}
\end{ex}
\vspace{0.5mm}

\begin{ex} (0,75 điểm)
Một cửa hàng sách có hai khu sách mới và sách cũ, mỗi khu được bán đồng giá. Bạn Mai chi $112\,500\text{ đồng}$ để mua $3$ cuốn sách mới và $4$ cuốn sách cũ. Bạn Linh chi $157\,500\text{ đồng}$ để mua $3$ cuốn sách mới và $10$ cuốn sách cũ. Tính giá mỗi cuốn sách mới và giá mỗi cuốn sách cũ.
\end{ex}
\vspace{0.5mm}

\begin{ex} (0,75 điểm)
Các tia nắng mặt trời tạo với mặt đất một góc xấp xỉ bằng $34^\circ$ và bóng của một ngọn tháp trên mặt đất dài $86\text{ m}$. Tính chiều cao của tháp (làm tròn đến hàng đơn vị mét).
\begin{center}
\begin{tikzpicture}[scale=1.1, line join=round, line cap=round, >=stealth]
  \coordinate (B) at (0,0); % Chan thap
  \coordinate (A) at (0,2.8); % Dinh thap
  \coordinate (C) at (4.2,0); % Dau bong thap tren mat dat
  \draw[thick] (-0.5,0) -- (5,0) node[below] {Mặt đất};
  \draw[line width=1.5pt, color=blue!70!black] (B) -- (A) node[midway, left] {Tháp $h$};
  \draw[thick, dashed] (B) -- (C) node[midway, below] {Bóng $86\text{ m}$};
  \draw[thick, color=orange!80!black] (A) -- (C) node[midway, above right] {Tia nắng};
  \draw (0,0.3) -- (0.3,0.3) -- (0.3,0);
  \draw[->, thick, orange!90!black] (3.2,0) arc (180:146:1);
  \node at (3.5,0.3) [orange!90!black] {$34^\circ$};
  \filldraw[black] (B) circle (1.5pt) node[below left] {$B$};
  \filldraw[black] (A) circle (1.5pt) node[above] {$A$};
  \filldraw[black] (C) circle (1.5pt) node[below right] {$C$};
\end{tikzpicture}
\end{center}
\end{ex}
\vspace{0.5mm}

\begin{ex} (0,5 điểm)
Cho tam giác $ABC$ có $\widehat{A} = 90^\circ$. Chứng minh rằng $\sin B + \cos C = 2\sin B$.
\end{ex}
\vspace{0.5mm}


\enlargethispage{1.5cm}
\begin{center}
\rule{4cm}{0.4pt}\\
\textit{--- HẾT ---}
\end{center}
\end{document}
